\section{Introduction}
%
%context & need
Amphibians are currently not explicitly considered in the environmental risk assessment (ERA) of chemicals. 
It is rather assumed that assessments based on aquatic vertebrates (e.g. fish) 
and terrestric vertebrates (birds and mammals) are also protective for amphibians. 
While amphibian larvae are in fact often less sensitive to chemicals 
than amphibian aquatic stages \citep{weltjeCOMPARATIVEACUTECHRONIC2013}, 
this is not the case for all chemicals \citep{glabermanEvaluatingRoleFish2019}, 
and the process of metamorphosis represents a vulnverable point in the life-hisory of amphibians, 
which is not covered by using other vertebrates as proxies.\\
Furthermore, environmental exposure occurrs in mixtures. 
The relative effect of mixtures within the metamorphosis phase is so far largely unexplored, 
and it cannot be expected that mixture toxicity data will be generated for a wide range 
of species and chemicals, necessitating the development of predictive models.\\
%
% goal
In this study, we demonstrate the possible use of toxicokinetic-toxicodynamic (TKTD) models 
to support a reliable ERA of chemicals for amphibians, including the extrapolation 
of chemical effects from larvae to metamorphs and the prediction of mixture effects. 
In general, this approach aims to maximize the amount of knowledge gained 
from experimental data, 
and to minimize the need for future animal testing. \\
%
\subsection{Setting the scene: predictive modelling of chemical toxicity with Dynamic Energy Budgets}
%
Dynamic Energy Budget (DEB) models are by now the most established models 
for predictive and interpretable modelling of chemical toxicity. 
These models describe the individual as a mass balance, 
starting with the amount of ingested food and ending with the central metabolic processes determining an 
individual's life-history:  
maintenance, growth,: maturation and reproduction \citep{kooijmanDynamicEnergyBudget2010,jagerDEBkissQuestSimplest2013}.
DEB models can be extended to account for the idiosyncracies of the amphibian life cycle \citep{romoliDynamicEnergyBudget2024,hansulDynamicModellingAmphibian}.
Effects of chemicals can be incorporated into DEB models by coupling them to a toxicokinetic-toxicodynamic (TKTD) submodel. 
These translate the external chemical exposure to an effect on the energy budgets, 
e.g. a decrease in a conversion efficiency, increase in maintenance costs or change in resource allocation.
\textit{A priori}, there is no good reason to deviate from 
the standard approach for the application to amphibians. 
In general, it can be necessary to include other PMoAs to explain 
observed effects, 
or to consider the simultaneous action of multiple PMoAs \citep{kochMakingSenseLifeHistory2021}.
% 
\subsection{Objective}
%
The objective of the present study was to test the applicability of DEB-TKTD models to 
analyze toxicity data in amphibians, specicially the anuran \textit{Discoglossus galganoi}, 
in order to draw conclusions about mechanisms of toxicity and synthesize effects through different exposure patterns, 
on different endpoints and on different life stages. \\ 
We briefly describe and motivate the model structure, and report results on three case studies in \textit{D. galganoi}, 
concerning the estimation of DEB parameters from control data with account for biological variabiltiy,
the comparative effects of the herbicide 2,4 Dichlorphenoxyacetic acid (2,4D) and the insecticide product Sivanto,
and the 
%The objective of the present study is to test the applicability 
%of DEB-TKTD models, specicially the \textit{AmphiDEB} model, 
%to the predictive modelling of combined effects of chemicals and pathogens.\\
%We divide \textit{applicability} into two aspects: 
%Firstly, the potential to explain observed effects 
%on the basis of energy budgets (model calibration). \\
%Secondly, the potential to predict combined effects 
%of chemicals and pathogens when the model is only calibrated 
%to single-stressor effects (model validaton).\\
%The approach of this study can be roughly divided into four steps: 
%\begin{enumerate}
%    \item Calibrate the model with data on chemical toxicity, assessing the potential of the model to explain observed chemical effects. 
%    \item Calibrate the model with data on pathogen loads and effects, assessing the potential of the model to explain observed pathogen effects.
%    \item Generate predictions of combined chemical and pathogen effects, assessing the predictive power of the model. 
%    \item If needed, re-calibrate the model on combined chemical and pathogen effects, in order understand discrepancies between predicted and observed effects revealed in the previous step.  
%\end{enumerate}  
